\documentclass[../../../include/open-logic-section]{subfiles}

\begin{document}

\olfileid{sth}{z}{sep}
\olsection{في الفصل}

نستفتح بمبدأ يقوم مقام الاستيعاب الساذج:

\begin{axiom}[مخطط الفصل] لكل صيغة~$\phi(x)$ يكون الحكم الآتي مسلّمة:
لكل~$A$ توجد المجموعة~$\Setabs{x \in A}{\phi(x)}$.
\end{axiom}

وليس هذا مسلّمة مفردة، بل \emph{مخطط} مسلّمات؛ فمسلّمات الفصل
\emph{غير متناهية العدد}، لكل صيغة~$\phi(x)$ واحدة. وصيغته المكافئة
المعهودة هي:

\begin{defish}
لكل صيغة~$\phi(x)$ لا يرد فيها «$S$»، يكون الآتي مسلّمة:
\[
	\forall A \exists S \forall x(x \in S \liff (\phi(x) \land x \in A)).
\]
\end{defish}

وبحسب الاصطلاح المقرر في صدر~\olref[sth][][]{part}، يجوز أن تشتمل
الصيغ~$\phi$ في مسلّمات الفصل على معاملات.\footnote{لشرح معنى ذلك،
  انظر المناقشة التي تلي مباشرة
  \olref[sfr][infinite][induction]{natinductionschema}.}

والتصور التراكمي-التكراري يسوغ الفصل من فوره. خذ مجموعة~$A$؛ فإن~$A$
قد تكونت عند مرحلة~$S$ ما بموجب~\stageshier. ولأن~$A$ تكونت في
المرحلة~$S$، تكونت عناصر~$A$ كلها قبل المرحلة~$S$ بموجب~\stagesacc.
فالمجموعات التي تنتمي إلى~$A$ وتحقق~$\phi$ تكونت أيضًا قبل
المرحلة~$S$، ومن ثم تُجمع في~$\Setabs{x \in A}{\phi(x)}$ عند
المرحلة~$S$ بموجب~\stagesacc.

وهذا، على خلاف الاستيعاب الساذج، سالم من مفارقة راسل؛ فلا يثبت به
وجود~$\Setabs{x}{x \notin x}$ مطلقًا. وإنما إذا \emph{أُعطينا}
مجموعة~$A$ أثبتنا~$R_A = \Setabs{x \in A}{x \notin x}$. ولا يلزم
من ذلك إلا~$R_A \notin R_A$ و$R_A \notin A$، وليس فيهما محذور.

لكن للفصل نتيجة قريبة جديرة بالتنبيه:

\begin{thm}\ollabel{thm:NoUniversalSet}
يمتنع وجود مجموعة \emph{شاملة}؛ أي لا توجد~$\Setabs{x}{x = x}$.
\end{thm}

\begin{proof}
افرض بالخلف أن~$V$ مجموعة شاملة. فعندئذ توجد، بموجب الفصل،
$R = \Setabs{x \in V}{x \notin x} = \Setabs{x}{x \notin x}$،
وهذا يناقض مفارقة راسل.
\end{proof}

وعدم المجموعة الشاملة، أو انفتاح هرم المجموعات، من أصول التصور
التراكمي-التكراري. ويمكن إدراك النتيجة من وجه آخر: فلو كانت مجموعة
شاملة لكانت !!a{element} في نفسها. لكن كل مجموعة، بحسب التصور، لا تظهر
أول مرة في الهرم إلا في أول مرحلة تلي جميع ما ينتمي إليها من
!!{element}. فيلزم أن \emph{لا} تنتمي مجموعة إلى نفسها؛ وإلا لوجب أن
يكون أول ظهورها بعد المرحلة التي وقع فيها أول ظهورها، وهو محال. (وسنرى
في
\olref[sth][spine][rank]{defnsetrank} كيف نجعل هذا الشرح أدق باستعمال
مفهوم «رتبة» المجموعة. لكننا سنحتاج إلى وضع بضع مسلّمات أخرى قبل ذلك.)
% تصحيح ترقيمي (0539-N1): أعيد قوس البدء للملاحظة الاعتراضية.

ومن الفصل والامتدادية تلزم أيضًا النتائج الآتية.

\begin{prop}\ollabel{prop:emptyexists}
متى وُجدت مجموعة ما وُجدت~$\emptyset$.
\end{prop}

\begin{proof}
  إذا كانت~$A$ مجموعة، وُجدت
  $\emptyset = \Setabs{x \in A}{x \neq x}$ بموجب الفصل.
\end{proof}

\begin{prop}
توجد~$A \setminus B$ متى أُعطيت مجموعتان~$A$ و$B$.
\end{prop}

\begin{proof}
توجد~$A \setminus B = \Setabs{x \in A}{x \notin B}$ بموجب الفصل.
\end{proof}

ويكاد الفصل يثبت وجود التقاطعات الاعتباطية أيضًا:

\begin{prop}\ollabel{prop:intersectionsexist}
  متى كان~$A \neq \emptyset$، وُجدت
  $\bigcap A = \Setabs{x}{(\forall y \in A)x \in y}$.
\end{prop}

\begin{proof}
افرض~$A \neq \emptyset$؛ فيوجد~$c \in A$. وعندئذ
$\bigcap A = \Setabs{x}{(\forall y \in A)x \in y} =
\Setabs{x \in c}{(\forall y \in A)x \in y}$، وهي موجودة بموجب الفصل.
\end{proof}

ولا بد من الشرط~$A \neq \emptyset$؛ لأن~$\bigcap \emptyset$، لو وُجد،
لكان المجموعة الشاملة بصدق فارغ، وهذا مخالف لـ
\olref{thm:NoUniversalSet}.

\end{document}
